%% title: Lubukusu logophors and perspectival agreement
%% category: phenomena

\citet{diercks2013} analyses complementizer agreement in Lubukusu, where sometimes the complementizer of an embedded clause agrees in noun class with the closest higher subject. The pattern is presented in \Next, where glossed numbers refer to noun classes:

\ex \begingl \gla Ba-ba-ndu ba-bol-el-a Alfredi \textbf{ba-li} a-kha-khil-e. //
    \glb 2-2-people 2s-said-\textsc{ap-fv} 1Alfred 2-that 1s-\textsc{fut}-conquer //
    \glft `The people told Alfred that he will win.' //
    \endgl
\xe

Here the embedded complementizer \textit{bali} is agreeing with the matrix subject \textit{babandu}.

This would provide clear evidence for an Upwards Agree that can make reference to subject positions specifically; however, Diercks argues that in fact there is no Upwards Agree going on. Instead, the subject whose features are agreed with binds a null anaphor in the embedded clause, which lies in the specifier of the complementizer, agreeing with it.

Diercks presents some semantic evidence, evidence from blocking, and parallels with logophors, to show that the properties of the dependency in \Last{} are not those one would expect from Upwards Agreement.

This paper is important for my thesis, as it details one of the most difficult agreement cases to explain in derivational frameworks, and provides a pretty good account for it.
